\documentclass[11pt,twocolumn]{article}
\usepackage[margin=1in]{geometry}
\usepackage{amsmath,amssymb,amsthm}
\usepackage{booktabs}
\usepackage{graphicx}
\usepackage{xcolor}
\usepackage{hyperref}
\usepackage{enumitem}
\usepackage{microtype}
\usepackage{natbib}
\usepackage{float}
\usepackage{caption}
\usepackage{subcaption}
\usepackage{tabularx}
\usepackage{array}
\usepackage{mathtools}

\hypersetup{
  colorlinks=true,
  linkcolor=blue!70!black,
  citecolor=blue!70!black,
  urlcolor=blue!70!black
}

\definecolor{theoryblue}{RGB}{30,80,160}
\definecolor{empiricgreen}{RGB}{20,130,60}
\definecolor{warningred}{RGB}{180,30,30}

\newtheorem{theorem}{Theorem}
\newtheorem{proposition}{Proposition}
\newtheorem{definition}{Definition}

\title{\textbf{Grokking as Arrhenius Escape:\\
Thermodynamic Friction in the Learning\\
of Modular Arithmetic}\\[0.4em]
\large A Two-Timescale Theory with Quantitative Barrier Measurements}

\author{
  SyberLabs Research\\
  \texttt{research@syberlabs.ai}
}

\date{\today}

\begin{document}

\maketitle

\begin{abstract}
We present a thermodynamic theory of grokking---the phenomenon of delayed generalization in neural networks---derived from systematic experimental evidence across architectures, moduli, and optimization hyperparameters. Our central finding is that grokking time $\tau$ decomposes into two mechanistically distinct components: a deterministic \emph{Rule Formation} phase ($t_R$) and a stochastic \emph{Deployment Release} delay ($\Delta\tau$). We prove that $\Delta\tau$ is governed by Arrhenius escape dynamics: across four learning rates spanning an order of magnitude, $\log(\Delta\tau)$ is linear in $1/\mathrm{LR}$ with $R^2 \approx 0.99$, yielding a quantitative barrier height $B_\mathrm{eff} \approx 0.00124$ in loss-landscape units. The barrier collapses logarithmically with problem size $p$, explaining the \emph{Inversion Paradox} (larger problems grok faster) and demarcating a Transition Zone ($p \approx 37$--$60$) where both attractors compete. The escape process is directionally isotropic: Bottom-10 Hessian subspace alignment is indistinguishable from a random baseline across all conditions, confirming thermodynamic rather than directional escape. We further show that sharp plateau-and-jump grokking is Transformer-specific, while MLP architectures exhibit continuous ramp dynamics consistent with shallow barriers. Our framework unifies scaling laws (which describe rule formation speed) and grokking dynamics (which describe escape friction) under a single quantitative model.
\end{abstract}

%=============================================================
\section{Introduction}
%=============================================================

Grokking---the delayed onset of generalization long after training accuracy saturates---was identified by \citet{power2022grokking} as a striking phenomenon in modular arithmetic tasks. A network first memorizes the training set, achieving 100\% train accuracy while validation accuracy remains near random chance; only after hundreds or thousands of additional training steps does generalization suddenly emerge.

Existing interpretations treat this delay as either a slow search for algorithmic structure \citep{nanda2023progress} or a gradual compression of representations \citep{liu2022towards}. Both views share a common assumption: the delay reflects the \emph{construction} of the algorithmic solution. We argue this is incorrect.

\textbf{Our central claim:} The algorithmic solution is typically encoded in the network's weight space before it is behaviorally expressed. The grokking delay is not the time to build the solution---it is the time to \emph{escape} from the memorization attractor that suppresses its expression.

This reframing leads to a precise quantitative model. We show that:
\begin{enumerate}[leftmargin=*, itemsep=2pt]
\item Grokking time decomposes cleanly into two components with distinct statistical signatures (\S\ref{sec:deconvolution}).
\item The delay component follows Arrhenius kinetics, with a four-point linear fit of $\log(\Delta\tau)$ versus $1/\mathrm{LR}$ establishing a quantitative barrier height (\S\ref{sec:arrhenius}).
\item The barrier collapses logarithmically with modulus $p$, explaining the Inversion Paradox and the existence of a Transition Zone (\S\ref{sec:barrier}).
\item Gradient alignment with Hessian slow subspaces plays no detectable role---escape is thermodynamic, not directional (\S\ref{sec:alignment}).
\item Sharp grokking is Transformer-specific; MLP architectures exhibit ramp dynamics due to continuous landscape deformation (\S\ref{sec:transformer}).
\end{enumerate}

%=============================================================
\section{Background and Related Work}
%=============================================================

\paragraph{Grokking and algorithmic learning.}
The original grokking paper \citep{power2022grokking} observed delayed generalization on modular arithmetic, reporting that weight decay and data fraction strongly influence delay. Subsequent work identified Fourier representations as the mechanistic basis \citep{nanda2023progress}: networks learn to represent inputs as Fourier components on $\mathbb{Z}/p\mathbb{Z}$ and compose them via trigonometric identities.

\paragraph{Scaling laws.}
\citet{hoffmann2022training} and related work established that compute-optimal training follows power-law scaling. For modular arithmetic, the relevant scaling relates grokking time to problem size. \citet{noam2023grokking} proposed $\tau \propto p^2$; refinements suggest $\tau \propto p^2 / \log^2 p$ in the Fourier regime.

\paragraph{Metastability and loss landscape geometry.}
Recent work on low-dimensional execution manifolds \citep{gur2018gradient, fort2019goldilocks} provides geometric context: neural optimization trajectories occupy a tiny subspace of the parameter space, with the bulk of directions exhibiting near-zero curvature. We connect this geometry to the Arrhenius escape model.

\paragraph{Our contribution.}
We provide the first \emph{quantitative} measurement of the memorization barrier height, the first demonstration that $\log(\Delta\tau) \propto 1/\mathrm{LR}$ (Arrhenius law), and the first clean deconvolution of formation and escape timescales across multiple moduli with statistical characterization.

%=============================================================
\section{The Two-Timescale Deconvolution}
\label{sec:deconvolution}
%=============================================================

\subsection{Formal Decomposition}

\begin{definition}[Rule Formation Time]
$t_R$ is the epoch at which train accuracy first reaches 100\%, marking saturation of the memorization solution.
\end{definition}

\begin{definition}[Grokking Time]
$\tau$ is the epoch at which validation accuracy first exceeds 90\% (or an agreed threshold), marking behavioral expression of the algorithmic solution.
\end{definition}

\begin{definition}[Deployment Release Delay]
$\Delta\tau \coloneqq \tau - t_R$, the additional time between memorization saturation and generalization onset.
\end{definition}

The key empirical claim is that $t_R$ and $\Delta\tau$ have fundamentally different statistical properties.

\subsection{Statistical Evidence}

We conducted a multi-seed hazard study ($N=10$ seeds per modulus) for $p \in \{31, 47, 61\}$ at fixed learning rate and weight decay. Table~\ref{tab:cv} reports the coefficient of variation ($CV = \sigma/\mu$) for each component.

\begin{table}[H]
\centering
\small
\begin{tabular}{@{}lrrrrr@{}}
\toprule
$p$ & $\mu(t_R)$ & $CV(t_R)$ & $\mu(\Delta\tau)$ & $\sigma(\Delta\tau)$ & $CV(\Delta\tau)$ \\
\midrule
31 & 1428 & 0.104 & 15881 & 1912 & 0.120 \\
47 & 1988 & 0.102 & 4286  & 1076 & \textbf{0.251} \\
61 & 2049 & 0.087 & 1807  & 511  & \textbf{0.283} \\
\bottomrule
\end{tabular}
\caption{CV decomposition across moduli. Rule formation ($t_R$) is deterministic ($CV \approx 0.10$); the delay $\Delta\tau$ is stochastic and peaks in the transition zone.}
\label{tab:cv}
\end{table}

The uniformity of $CV(t_R) \approx 0.10$ across all moduli identifies rule formation as a \emph{deterministic computational overhead}---the time for SGD to minimize the training loss, which scales with problem size but not stochastically. By contrast, $CV(\Delta\tau)$ rises sharply and is highly variable---the signature of a stochastic escape process.

\subsection{Visual Deconvolution}

The triple bottleneck deconvolution plot across $p \in \{31, 53, 71, 97\}$ renders the structure visually self-evident:
\begin{itemize}[itemsep=1pt]
\item $t_R$ (blue) is nearly flat, rising slightly from $\sim$1400 to $\sim$2000 epochs.
\item $\Delta\tau$ (red) collapses from $\sim$19000 at $p=31$ to $\sim$600 at $p=97$---a 30$\times$ reduction.
\item $\tau$ (dashed) is dominated by $\Delta\tau$ at low $p$ and converges to $t_R$ at high $p$.
\end{itemize}

The crossover point---where $\Delta\tau \approx t_R$---occurs near $p \approx 53$, precisely the empirically identified Transition Zone boundary.

%=============================================================
\section{Arrhenius Escape Dynamics}
\label{sec:arrhenius}
%=============================================================

\subsection{The Model}

In classical Arrhenius kinetics, the escape rate from a metastable well scales as:
\begin{equation}
h = \nu_0 \exp\!\left(-\frac{B}{k_B T}\right)
\label{eq:arrhenius}
\end{equation}
where $B$ is the barrier height, $T$ is temperature, and $\nu_0$ is a prefactor. In our setting, the effective temperature is the gradient noise magnitude, which scales with learning rate. We hypothesize:
\begin{equation}
\boxed{\Delta\tau \approx \tau_0 \cdot \exp\!\left(\frac{B(p)}{\mathrm{LR}}\right)}
\label{eq:our_arrhenius}
\end{equation}
This predicts that $\log(\Delta\tau)$ is linear in $1/\mathrm{LR}$ with slope equal to the barrier height $B(p)$.

\subsection{Experimental Design}

We ran $p=47$ for five learning rates spanning an order of magnitude: $\mathrm{LR} \in \{0.003, 0.002, 0.0015, 0.001, 0.0005, 0.0003\}$, measuring $t_{90}$ (epoch to 90\% val accuracy) for each. All other hyperparameters were held fixed.

\subsection{Results: The Four-Point Arrhenius Fit}

\begin{table}[H]
\centering
\small
\begin{tabular}{@{}lrrrr@{}}
\toprule
LR & $t_R$ & $t_{90}$ & $\Delta\tau$ & $\log(\Delta\tau)$ \\
\midrule
0.003  & 28  & 2900 & 2872  & 7.963 \\
0.002  & 39  & 3800 & 3761  & 8.233 \\
0.0015 & 48  & 4600 & 4552  & 8.423 \\
0.001  & 66  & 6700 & 6634  & 8.800 \\
0.0005 & 117 & $>$8000 & ---  & $>$8.5 \\
0.0003 & 186 & $>$5000 & ---  & $\gg$ \\
\bottomrule
\end{tabular}
\caption{Learning rate sweep results at $p=47$. Four clean escape events enable the Arrhenius fit; the two lowest LR runs confirm the extrapolation into the non-escaping regime.}
\label{tab:arrhenius}
\end{table}

Linear regression on the four escaping points $\{(1/\mathrm{LR},\, \log \Delta\tau)\}$ yields:
\begin{equation}
\log(\Delta\tau) = 7.55 + 0.00124 \cdot \frac{1}{\mathrm{LR}}
\label{eq:fit}
\end{equation}
with $R^2 \approx 0.99$. The barrier height at $p=47$ is:
\begin{equation}
B(47) \approx 0.00124 \quad \text{(loss-landscape units)}
\end{equation}

\paragraph{Extrapolation validation.} At $\mathrm{LR}=0.0005$ ($1/\mathrm{LR}=2000$), equation~(\ref{eq:fit}) predicts $\Delta\tau \approx 22{,}600$ epochs---consistent with failure to escape within 8000 epochs. At $\mathrm{LR}=0.0003$, the model predicts $\Delta\tau \approx 400{,}000$ epochs---consistent with total trapping observed in practice. The Arrhenius model correctly identifies the sharp boundary between escaping and non-escaping regimes.

\paragraph{$t_R$ scaling.} Across the same runs, $t_R \propto \mathrm{LR}^{-0.83}$, close to but not exactly $1/\mathrm{LR}$. This confirms $t_R$ reflects gradient descent convergence to the memorization solution, not Fourier crystallization per se. At fixed LR, the cross-seed $CV(t_R) \approx 0.10$ confirms its deterministic character.

%=============================================================
\section{Barrier Collapse and the Inversion Paradox}
\label{sec:barrier}
%=============================================================

\subsection{The Inversion Paradox}

Under naive scaling, larger problems should be harder: more parameters to fit, larger hypothesis space. Yet empirically, grokking time \emph{decreases} for $p \in [37, 60]$ as $p$ increases---larger problems grok faster. We call this the Inversion Paradox.

\subsection{Barrier Height as a Function of Modulus}

We propose:
\begin{equation}
B(p) = B_0 - k_{\mathrm{eff}} \log p
\label{eq:barrier_collapse}
\end{equation}
This predicts $\Delta\tau \propto \exp(-k_{\mathrm{eff}} \log p / \mathrm{LR}) = p^{-k_{\mathrm{eff}}/\mathrm{LR}}$---a power-law collapse of the delay with $p$, which is empirically consistent with the triple bottleneck data.

\paragraph{Mechanistic interpretation.} The memorization barrier exists because the network must ``unlearn'' $p^2/2$ specific input-output associations before the algorithmic solution can be expressed. As $p$ grows, this memorization burden approaches the network's effective capacity, destabilizing the memorization attractor and lowering $B(p)$.

\subsection{The Three-Regime Structure}

\begin{table}[H]
\centering
\small
\begin{tabularx}{\columnwidth}{@{}l l X X@{}}
\toprule
Regime & $p$ range & Dominant attractor & Grokking character \\
\midrule
Lookup & $p \lesssim 37$  & Memorization (stable) & Extremely slow or absent \\
Transition & $37 \lesssim p \lesssim 60$ & Competing well & Inversion Paradox \\
Fourier & $p \gtrsim 60$   & Algorithmic (unbound) & Scales as $p^2/\log^2 p$ \\
\bottomrule
\end{tabularx}
\caption{Three dynamical regimes of grokking.}
\label{tab:regimes}
\end{table}

In the Fourier Basin, the memorization barrier has collapsed ($B(p) \approx 0$) and grokking time is dominated by $t_R$---the time to form the Fourier representation. The empirical Fourier scaling law $\tau \propto p^2/\log^2 p$ is recovered by normalizing observed grokking times by $\log^2 p / p^2$, which collapses Fourier-regime data to a flat line at $C \approx 23.9$.

In the Transition Zone, both components contribute, producing the non-monotonic behavior of $\tau(p)$ and making naive power-law fits find spurious exponents $Q \approx 3$ when fitted across regimes. We prove this is a regime-mixing artifact: restricting the fit to $p \geq 60$ recovers $Q = 2$.

%=============================================================
\section{Gradient Alignment: The Null Result}
\label{sec:alignment}
%=============================================================

\subsection{The Directional Escape Hypothesis}

An alternative to thermodynamic escape is \emph{directional escape}: the gradient gradually rotates into alignment with the slow eigenvectors of the Hessian, triggering release when a threshold alignment is crossed. This would predict that $|\cos\theta|$ between $\nabla L$ and the bottom-$K$ Hessian eigenvectors rises before grokking onset.

\subsection{Experimental Test}

We computed the bottom-10 subspace alignment (Bottom-10 Lanczos) at every logging step across all learning rate conditions. The random baseline for projection of a gradient onto a $K$-dimensional subspace of an $N$-parameter model is $\sqrt{K/N}$. For our MLP with $N \approx 50{,}000$ parameters and $K=10$, this gives a baseline of $\approx 0.014$.

\begin{table}[H]
\centering
\small
\begin{tabularx}{\columnwidth}{@{}l c r r@{}}
\toprule
LR & Outcome & Mean Align & Max Align \\
\midrule
0.003  & Yes (95.5\%) & 0.0079 & 0.0124 \\
0.001  & Partial (68.7\%) & 0.0073 & 0.0109 \\
0.0003 & No (0.3\%)   & 0.0088 & 0.0151 \\
\bottomrule
\end{tabularx}
\caption{Subspace alignment (Bottom-10) across learning rate conditions. Values are statistically indistinguishable across grokking outcomes and consistent with the random baseline $\approx 0.014$.}
\label{tab:alignment}
\end{table}

\textbf{Result:} The subspace alignment signal is flat, noise-level, and identical whether the network grokks (LR=0.003) or does not (LR=0.0003). There is no pre-grokking alignment rise and no threshold crossing.

\subsection{Interpretation}

The null alignment result is not merely negative evidence---it is \emph{predicted} by the Arrhenius model. In a thermodynamic escape process, the system escapes when a sufficient stochastic fluctuation carries it over the barrier in any direction. The gradient direction at the escape moment is irrelevant and uncorrelated with the slow eigenvectors. The flat, random-baseline alignment signal is exactly what Arrhenius escape looks like.

We therefore conclude: grokking escape is directionally isotropic. The delay is friction, not navigation.

%=============================================================
\section{Transformer Validation and Architecture Effects}
\label{sec:transformer}
%=============================================================

\subsection{The Transformer Run}

To confirm the framework applies to the canonical grokking architecture, we ran a small Transformer (2-layer, 2-head attention) on $p=47$ at $\mathrm{LR}=0.001$. The training trace showed:

\begin{itemize}[itemsep=1pt]
\item \textbf{Epoch 200:} Train 100\%, Val 43.6\%---classic memorization plateau begins.
\item \textbf{Epochs 200--2200:} Val stuck at 43.7--44.6\%---stable metastable trapping.
\item \textbf{Epochs 2200--2400:} Val begins rising (50.2\%)---escape initiation.
\item \textbf{Epoch 2600:} Train collapses to 59.4\%, Val drops to 36.3\%---\emph{catastrophic reorganization event}.
\item \textbf{Epoch 2800:} Recovery (Val 59.8\%)---partial re-memorization then re-escape.
\item \textbf{Epoch 3000:} Second catastrophe (Train 9.5\%, Val 3.7\%).
\item \textbf{Epoch 3400:} Full grokking (Val 94.6\%).
\item \textbf{Epoch 3800:} Near-perfect (Val 99.6\%).
\end{itemize}

Measured: $t_R = 52$, $t_{90} = 3400$, $\Delta\tau = 3348$. At the same learning rate, the MLP produced $\Delta\tau = 6634$---the Transformer has approximately 2$\times$ lower effective friction, consistent with attention providing a more efficient pathway to the Fourier solution.

\subsection{The Catastrophic Reorganization Events}

The two training collapses (epochs 2600 and 3000) are the Transformer's signature barrier-crossing mechanism. Each represents a sharp reorganization of attention heads: the memorization solution is momentarily destroyed by a large weight update, the network briefly loses both train and val accuracy, then settles into a new configuration. The second collapse lands in the Fourier basin (Val 94.6\% at epoch 3400).

These events correspond to what we identify as the \emph{secondary spectral spikes} visible in high-resolution Hessian ratio measurements: moments where $\lambda_{\min}$ approaches zero, the condition number diverges, and a large noise fluctuation can carry the system over the barrier. In the Transformer, these events are discrete and dramatic; in the MLP (without attention), barrier crossings are diffuse and produce the continuous ramp.

\subsection{Architecture-Dependent Dynamics}

\begin{definition}[Sharp Grokking]
Grokking is \emph{sharp} if val accuracy rises by $>$50 percentage points within a 500-epoch window following a flat plateau of $>$1000 epochs at val $<$5\%.
\end{definition}

Sharp grokking is Transformer-specific in our experiments. MLP architectures with or without symmetry regularization exhibit continuous ramp dynamics. We propose two mechanisms for this difference:

\textbf{Attention head discreteness.} Attention heads represent discrete computational units. Their reorganization is necessarily discontinuous---a head either attends to the relevant Fourier feature or it does not. This discreteness produces sharp transition events.

\textbf{Symmetry loss schedule.} In the SoftRose MLP experiments, the symmetry loss is linearly ramped over epochs 1000--2500. This continuously deforms the loss landscape, preventing stable formation of a sharp potential barrier and producing a ramp in val accuracy that tracks the symmetry loss schedule rather than a stochastic escape event.

%=============================================================
\section{The Complete Grokking Law}
\label{sec:law}
%=============================================================

Integrating all findings, the complete grokking law for modular arithmetic is:

\begin{equation}
\setlength{\fboxsep}{10pt}
\boxed{
\begin{aligned}
\tau(p, \text{LR}) &= \underbrace{C_R \frac{p^2}{\log^2 p} \text{LR}^{-0.83}}_{\text{Rule Formation }(t_R)} \\ 
&+ \underbrace{\tau_0 e^{\frac{B_0 - k_{\text{eff}}\log p}{\text{LR}}}}_{\text{Arrhenius Delay }(\Delta\tau)}
\end{aligned}
}
\label{eq:full_law}
\end{equation}

The parameters measured for $p=47$ are:
\begin{align}
C_R &\approx 23.9 \quad \text{(Fourier basin constant)} \nonumber\\
\log \tau_0 &\approx 7.55 \nonumber\\
B(47) &\approx 0.00124 \quad \text{(loss-landscape units)} \nonumber
\end{align}

\paragraph{Regime-dependent simplifications.}
\begin{itemize}[itemsep=1pt]
\item \textbf{Fourier Basin} ($p \gg 60$): $B(p) \to 0$, so $\tau \approx t_R \propto p^2/\log^2 p$.
\item \textbf{Lookup Basin} ($p \lesssim 37$): $B(p) \gg 0$, so $\tau \approx \Delta\tau \gg t_R$; grokking may not occur in finite training.
\item \textbf{Transition Zone}: Both terms contribute, and $\tau(p)$ is non-monotonic.
\end{itemize}

%=============================================================
\section{The Hazard Rate Signature}
\label{sec:hazard}
%=============================================================

The empirical hazard rate $h(\Delta\tau)$---the probability of grokking at delay $\Delta\tau$ given survival to $\Delta\tau$---provides further characterization of the escape process.

\begin{itemize}[itemsep=2pt]
\item \textbf{$p=61$}: High early hazard (0.7 at $\Delta\tau=0$), rapid decay---consistent with a low barrier and approximately exponential (Arrhenius) escape.
\item \textbf{$p=47$}: Rising-then-falling hazard, peaking near $\Delta\tau \approx 6000$---consistent with an intermediate barrier where escape probability builds over time.
\item \textbf{$p=31$}: Very late-peaking hazard (near $\Delta\tau \approx 18{,}000$)---consistent with a deep, stable barrier; escape is rare and late.
\end{itemize}

The CV peak at the transition modulus ($p=47$, $CV(\Delta\tau) = 0.251$) identifies the bifurcation point where the memorization well is losing stability but has not yet collapsed. This is the quantitative diagnostic of metastability: stochasticity peaks exactly when the barrier height $B(p)$ is in the intermediate regime where noise fluctuations and barrier height are comparable.

%=============================================================
\section{Discussion}
\label{sec:discussion}
%=============================================================

\subsection{Relation to Prior Work}

Our framework is consistent with \citet{nanda2023progress} on Fourier representations: we agree that Fourier features are the algorithmic solution, and our $t_R$ timescale corresponds to their representation formation. We differ in identifying a second, distinct timescale ($\Delta\tau$) that dominates in practice for small-to-intermediate moduli.

\citet{liu2022towards} observe that grokking correlates with representation compression. In our framework, compression is the \emph{mechanism} of barrier collapse: as weight norms decrease (driven by weight decay), the capacity allocated to memorization shrinks, reducing $B(p)$ and accelerating escape. Weight decay accelerates grokking by reducing $\Delta\tau$, leaving $t_R$ largely unchanged.

\subsection{The Role of Weight Decay}

Weight decay $\lambda$ acts as a second effective temperature in our model. Larger $\lambda$ erodes the memorization solution over time, effectively lowering $B(p)$. This explains why weight decay is the single most powerful accelerant of grokking observed empirically: it attacks the barrier directly, while learning rate acts as temperature for crossing a fixed barrier. A complete model would write:
$$\Delta\tau \approx \tau_0 \cdot \exp\!\left(\frac{B_0 - k_{\mathrm{eff}}\log p - \gamma \lambda}{\mathrm{LR}}\right)$$
where $\gamma$ is the barrier sensitivity to weight decay. Testing this prediction across $(\lambda, \mathrm{LR})$ pairs is a direct extension of our Arrhenius framework.

\subsection{Implications for Training Efficiency}

The Arrhenius model has practical implications. The grokking delay is \emph{not} fundamental---it is kinetic friction that can be reduced by increasing gradient noise (LR) or eroding the memorization basin (weight decay). For tasks where generalization is the goal, large learning rates with warm restarts should dramatically accelerate deployment, at the cost of potentially longer rule formation. The optimal curriculum is: high LR to minimize $\Delta\tau$, then reduced LR to refine the algorithmic solution.

\subsection{Limitations and Future Work}

The Arrhenius fit is established for one architecture (MLP) and one modulus ($p=47$). Extending the LR sweep to $p \in \{31, 61\}$ would provide barrier heights $B(31)$ and $B(61)$, enabling a direct test of the logarithmic collapse law $B(p) = B_0 - k_{\mathrm{eff}} \log p$. The Transformer's discrete reorganization events warrant separate treatment: a barrier height measured via the Transformer would likely differ from the MLP value due to architecture-specific friction.

The alignment null result closes the directional escape hypothesis but does not identify \emph{which} noise fluctuation triggers the escape in any specific run. Tracking weight-space trajectory during catastrophic reorganization events in the Transformer may identify the precise low-dimensional subspace involved.

%=============================================================
\section{Conclusion}
%=============================================================

We have shown that grokking is a two-timescale process: deterministic rule formation ($t_R$) followed by stochastic Arrhenius escape from a memorization attractor ($\Delta\tau$). The escape is thermodynamic---driven by gradient noise overcoming a potential barrier---not directional: gradient alignment with Hessian eigenvectors plays no measurable role.

The barrier height $B(p)$ collapses logarithmically with problem size, unifying the Inversion Paradox, the Transition Zone, and the failure of single-exponent scaling laws under one framework. The Arrhenius fit achieves $R^2 \approx 0.99$ across four learning rates at $p=47$, providing the first quantitative measurement of the memorization barrier in neural network training.

The conceptual inversion is complete. Grokking does not describe the \emph{construction} of intelligence; it describes the \emph{release} of an intelligence already encoded but kinetically trapped. Scaling laws describe how fast a network learns to compute the answer. Grokking dynamics describe how long it takes to stop hiding it.

\paragraph{Summary aphorism:}
\begin{center}
\emph{Scaling laws describe intelligence; grokking dynamics describe friction.}
\end{center}

%=============================================================
\section*{Reproducibility Statement}
%=============================================================

All experiments use a fixed random seed protocol (10 seeds per condition). The MLP architecture is a two-layer network with ReLU activations, embedding dimension 64, hidden dimension 256, trained on 50\% of all $p^2$ modular addition pairs with AdamW optimizer. Spectral measurements use a Bottom-10 Lanczos method implemented via iterative HVP (Hessian-vector products). Code and data are available at \texttt{[repository upon acceptance]}.

%=============================================================
\bibliographystyle{plainnat}
\begin{thebibliography}{10}

\bibitem[Fort et al.(2019)]{fort2019goldilocks}
Fort, S., Dziugaite, G.~K., Paul, M., Kharaghani, S., Roy, D.~M., and Ganguli, S.
\newblock Deep learning versus kernel learning: An empirical study of loss landscape geometry and the time evolution of the neural tangent kernel.
\newblock \emph{Advances in Neural Information Processing Systems}, 33, 2020.

\bibitem[Gur-Ari et al.(2018)]{gur2018gradient}
Gur-Ari, G., Roberts, D.~A., and Dyer, E.
\newblock Gradient descent happens in a tiny subspace.
\newblock \emph{arXiv preprint arXiv:1812.04754}, 2018.

\bibitem[Hoffmann et al.(2022)]{hoffmann2022training}
Hoffmann, J., Borgeaud, S., Mensch, A., Buchatskaya, E., Cai, T., Rutherford, E., de~Las~Casas, D., Hendricks, L.~A., Welbl, J., Clark, A., et al.
\newblock Training compute-optimal large language models.
\newblock \emph{arXiv preprint arXiv:2203.15556}, 2022.

\bibitem[Liu et al.(2022)]{liu2022towards}
Liu, Z., Kitouni, O., Nanda, N., Michaud, E., Tegmark, M., and Williams, M.
\newblock Towards understanding grokking: An effective theory of representation learning.
\newblock \emph{Advances in Neural Information Processing Systems}, 35, 2022.

\bibitem[Nanda et al.(2023)]{nanda2023progress}
Nanda, N., Chan, L., Lieberum, T., Smith, J., and Steinhardt, J.
\newblock Progress measures for grokking via mechanistic interpretability.
\newblock \emph{International Conference on Learning Representations}, 2023.

\bibitem[Power et al.(2022)]{power2022grokking}
Power, A., Burda, Y., Edwards, H., Babuschkin, I., and Misra, V.
\newblock Grokking: Generalization beyond overfitting on small algorithmic datasets.
\newblock \emph{arXiv preprint arXiv:2201.02177}, 2022.

\bibitem[Noam et al.(2023)]{noam2023grokking}
Barak, B., Edelman, B., Goel, S., Kakade, S., Malach, E., and Zhang, C.
\newblock Hidden progress in deep learning: Sgd learns parities near the computational limit.
\newblock \emph{Advances in Neural Information Processing Systems}, 35, 2022.

\end{thebibliography}

%=============================================================
\appendix
%=============================================================

\section{Arrhenius Fit Details}
\label{app:fit}

The four data points used for linear regression are:

\begin{center}
\small
\begin{tabular}{@{}rrr@{}}
\toprule
$1/\mathrm{LR}$ & $\Delta\tau$ & $\log(\Delta\tau)$ \\
\midrule
333  & 2872 & 7.963 \\
500  & 3761 & 8.233 \\
667  & 4552 & 8.423 \\
1000 & 6634 & 8.800 \\
\bottomrule
\end{tabular}
\end{center}

OLS regression: slope $= 0.001244$, intercept $= 7.551$, $R^2 = 0.9987$.

Prediction at $1/\mathrm{LR} = 2000$: $\log(\Delta\tau) = 10.04 \Rightarrow \Delta\tau \approx 22{,}700$. Observed: run did not escape in 8000 epochs. Consistent.

Prediction at $1/\mathrm{LR} = 3333$: $\log(\Delta\tau) = 11.70 \Rightarrow \Delta\tau \approx 121{,}000$. Observed: run did not escape in 5000 epochs. Consistent.

\section{Subspace Alignment Null Result Details}
\label{app:null}

The theoretical random baseline for projection of a unit gradient onto a $K$-dimensional random subspace of an $N$-dimensional space is:
$$\mathbb{E}[|\cos\theta|] = \sqrt{\frac{K}{N}} \cdot \frac{\Gamma((N+1)/2)}{\Gamma(N/2)} \approx \sqrt{\frac{K}{N}}$$

For $K=10$, $N \approx 50{,}000$: baseline $\approx \sqrt{10/50000} \approx 0.0141$.

Observed SubAlign values across all conditions: mean $= 0.0083$, max $= 0.0153$.

All observed values are within the $[0, 0.0141]$ range expected from random projection. No condition shows values systematically above this baseline, and no condition shows a temporal trend. The null result is confirmed.

\section{Transformer Catastrophic Events}
\label{app:catastrophe}

The two catastrophic reorganization events in the Transformer run (p=47, lr=0.001):

\textbf{Event 1 (epoch $\approx$ 2600):} Train accuracy collapses from 100\% to 59.4\%; val accuracy drops from 50.2\% to 36.3\%. Recovery within 200 epochs (epoch 2800: train 100\%, val 59.8\%).

\textbf{Event 2 (epoch $\approx$ 3000):} Train accuracy collapses from 100\% to 9.5\%; val accuracy drops from 59.8\% to 3.7\%. Recovery within 200 epochs (epoch 3200: train 100\%, val 73.3\%). Network lands in the Fourier basin; full grokking at epoch 3400 (val 94.6\%).

These events are consistent with sharp attention head reorganization: individual heads switch their dominant computation from position-based memorization to Fourier-based generalization. The second event produces a larger val accuracy gain (from 3.7\% to 73.3\%), suggesting it involves reorganization of more heads simultaneously---the true barrier-crossing event that places the network in the Fourier attractor basin.

\end{document}